\documentclass[aspectratio=169]{beamer}

% Theme and color scheme
\usetheme{default}
\usecolortheme{default}

% Remove navigation symbols
\setbeamertemplate{navigation symbols}{}

% Simple color scheme for printing
\definecolor{darkblue}{RGB}{25,25,112}
\setbeamercolor{title}{fg=darkblue}
\setbeamercolor{frametitle}{fg=darkblue}
\setbeamercolor{structure}{fg=darkblue}

% Simple footline
\setbeamertemplate{footline}[frame number]

% Packages
\usepackage{amsmath}
\usepackage{amsfonts}
\usepackage{amssymb}
\usepackage{graphicx}
\usepackage{tikz}
\usepackage{booktabs}
\usepackage{array}
\usepackage{listings}
\usepackage{xcolor}

% Set graphics path for images (relative to slides directory)
\graphicspath{{../Images/}}

% Title information
\title[Chapter 3.12: Von Neumann Architecture]{Structured Computer Architecture\\Chapter 3.12: Von Neumann Architecture\\Fourth-Order Systems (4-OS)}
\author{Based on Structured Computer Architecture: An Order-Based Approach}
\date{}

% Custom command for highlighting concepts
\newcommand{\concept}[1]{\textbf{\color{darkblue}#1}}

\begin{document}

% Title slide
\begin{frame}
\titlepage
\end{frame}

% Table of contents
\begin{frame}{Outline}
\tableofcontents
\end{frame}

% Learning objectives
\begin{frame}{Learning Objectives}
By the end of this chapter, you will understand:
\begin{itemize}
\item The historical context and theoretical foundations of computation
\item The mathematical model of the Turing Machine
\item The decision problem (Entscheidungsproblem) and its significance
\item Von Neumann architecture principles and components
\item The stored program concept and its implications
\item Memory organization and addressing schemes
\item Instruction set architecture fundamentals
\item The relationship between hardware and software
\end{itemize}
\end{frame}

% Section 1: Historical Context and Theoretical Foundations
\section{Historical Context and Theoretical Foundations}

\begin{frame}{The Turning Point: Fourth-Order Systems}
\begin{columns}
\begin{column}{0.6\textwidth}
\concept{Fourth Loop Significance}:
\begin{itemize}
\item Meeting point of physical and informational structures
\item Circuit-information combination determines functionality
\item Architecture as unifying concept
\item Synthesis of symbolic and physical structures
\end{itemize}

\vspace{0.3cm}
\concept{Information as Control}:
\begin{itemize}
\item Information takes control at this level
\item Programs as symbolic structures
\item Mathematical models enable computation
\item Digital systems gain autonomy through loops
\end{itemize}
\end{column}
\begin{column}{0.4\textwidth}
\begin{center}
\textbf{Architecture Concept}
\vspace{0.3cm}

\tikz{
\node[draw, rectangle, minimum width=2.5cm, minimum height=1cm] (phys) at (0,2) {Physical Structure};
\node[draw, rectangle, minimum width=2.5cm, minimum height=1cm] (info) at (0,0) {Information Structure};
\node[draw, rectangle, minimum width=2.5cm, minimum height=0.8cm, fill=blue!20] (arch) at (0,1) {Architecture};
\draw[<->] (phys) -- (arch);
\draw[<->] (arch) -- (info);
}
\end{center}
\end{column}
\end{columns}
\end{frame}

\begin{frame}{Origins of Computation Theory}
\begin{columns}
\begin{column}{0.5\textwidth}
\concept{The Epimenides Paradox} (7th-6th century BC):
\begin{itemize}
\item "All Cretans are liars" - spoken by a Cretan
\item First recorded decision problem
\item Cannot determine truth or falsehood
\item Foundation of logical paradoxes
\end{itemize}

\vspace{0.3cm}
\concept{Historical Timeline}:
\begin{itemize}
\item Ancient Greece: Logical paradoxes
\item 1900: Hilbert's 23 problems
\item 1928: Entscheidungsproblem formulated
\item 1931: Gödel's incompleteness theorems
\item 1936: Turing machine concept
\end{itemize}
\end{column}
\begin{column}{0.5\textwidth}
\begin{center}
\textbf{Decision Problem Evolution}
\vspace{0.3cm}

\tikz{
\node[draw, circle, minimum width=1.5cm] (epi) at (0,3) {Epimenides};
\node[draw, circle, minimum width=1.5cm] (hil) at (0,2) {Hilbert};
\node[draw, circle, minimum width=1.5cm] (god) at (0,1) {Gödel};
\node[draw, circle, minimum width=1.5cm] (tur) at (0,0) {Turing};
\draw[->] (epi) -- (hil);
\draw[->] (hil) -- (god);
\draw[->] (god) -- (tur);
}
\end{center}
\end{column}
\end{columns}
\end{frame}

% Section 2: The Decision Problem
\section{The Decision Problem}

\begin{frame}{Hilbert's Entscheidungsproblem (1928)}
\begin{columns}
\begin{column}{0.6\textwidth}
\concept{The Question}:
\begin{quote}
\textit{"Is there an algorithm that takes a statement formulated in a formal system and answers with 'yes' or 'no' whether the statement is true or false in the given system?"}
\end{quote}

\vspace{0.3cm}
\concept{Significance}:
\begin{itemize}
\item Fundamental question about computability
\item Seeks universal decision procedure
\item Bridges logic and computation
\item Motivates formal system development
\end{itemize}
\end{column}
\begin{column}{0.4\textwidth}
\begin{center}
\textbf{Decision Process}
\vspace{0.3cm}

\tikz{
\node[draw, rectangle, minimum width=2cm, minimum height=0.8cm] (stmt) at (0,2) {Statement};
\node[draw, rectangle, minimum width=2cm, minimum height=0.8cm] (alg) at (0,1) {Algorithm};
\node[draw, rectangle, minimum width=2cm, minimum height=0.8cm] (ans) at (0,0) {Yes/No};
\draw[->] (stmt) -- (alg);
\draw[->] (alg) -- (ans);
}
\end{center}
\end{column}
\end{columns}
\end{frame}

\begin{frame}{Gödel's Revolutionary Answer (1931)}
\begin{columns}
\begin{column}{0.6\textwidth}
\concept{Gödel's Incompleteness Theorem}:
\begin{quote}
\textbf{\textit{"In a formal system one can correctly construct true structures whose truth cannot be proven or disproved"}}
\end{quote}

\vspace{0.3cm}
\concept{Implications}:
\begin{itemize}
\item No universal decision algorithm exists
\item Formal systems have inherent limitations
\item Truth vs. provability distinction
\item Foundation for computability theory
\end{itemize}

\vspace{0.3cm}
\concept{Impact on Computer Science}:
\begin{itemize}
\item Triggered four independent works (1936)
\item Church, Kleene, Post, Turing
\item Founded theoretical computer science
\end{itemize}
\end{column}
\begin{column}{0.4\textwidth}
\begin{center}
\textbf{Incompleteness Concept}
\vspace{0.3cm}

\tikz{
\draw[thick] (0,0) circle (2cm);
\node at (0,2.3) {Formal System};
\draw[thick] (0,0) circle (1.2cm);
\node at (0,1.5) {Provable};
\node at (0,0.5) {Statements};
\node[fill=red!30] at (1.5,0.5) {True but};
\node[fill=red!30] at (1.5,0.2) {Unprovable};
}
\end{center}
\end{column}
\end{columns}
\end{frame}

% Section 3: Turing Machine
\section{The Turing Machine}

\begin{frame}{Turing Machine Components}
\begin{columns}
\begin{column}{0.6\textwidth}
\concept{Three Essential Components}:
\begin{enumerate}
\item \textbf{Infinite Tape}: Stores symbols from finite alphabet
\item \textbf{Read/Write Head}: Accesses tape, moves left/right
\item \textbf{Finite Automaton}: Controls head based on state and symbol
\end{enumerate}

\vspace{0.3cm}
\concept{Operation Principle}:
\begin{itemize}
\item Each clock cycle: read symbol, update state
\item Head moves one position left or right
\item Write new symbol based on current state
\item Finite automaton determines all actions
\end{itemize}

\vspace{0.3cm}
\concept{Mathematical Abstraction}:
\begin{itemize}
\item Infinite tape handles unknown space requirements
\item Finite computation in practice
\item Universal model of computation
\end{itemize}
\end{column}
\begin{column}{0.4\textwidth}
\begin{center}
\textbf{Turing Machine Structure}
\vspace{0.3cm}

\tikz{
% Tape
\draw[thick] (-2,2) -- (2,2);
\draw[thick] (-2,1.5) -- (2,1.5);
\foreach \x in {-1.5,-1,-0.5,0,0.5,1,1.5} {
  \draw (\x,1.5) -- (\x,2);
}
\node at (0,1.75) {0};
\node at (0.5,1.75) {1};
\node at (-0.5,1.75) {1};

% Head
\draw[thick] (-0.2,1.5) -- (-0.2,1) -- (0.2,1) -- (0.2,1.5);
\draw[->] (0,1) -- (0,0.5);

% Finite Automaton
\draw[thick] (-0.5,0.5) rectangle (0.5,0);
\node at (0,0.25) {FA};
}
\end{center}
\end{column}
\end{columns}
\end{frame}

\begin{frame}{Digital Implementation of Turing Machine}
\begin{columns}
\begin{column}{0.6\textwidth}
\concept{Digital Components}:
\begin{itemize}
\item \textbf{Infinite RAM}: Replaces infinite tape
\item \textbf{Up-Down Counter}: Addresses memory locations
\item \textbf{Finite Automaton}: Same control logic
\end{itemize}

\vspace{0.3cm}
\concept{Practical Considerations}:
\begin{itemize}
\item Finite memory in real implementations
\item Counter provides sequential access
\item Memory addressing enables random access
\item Control logic remains unchanged
\end{itemize}

\vspace{0.3cm}
\concept{Theoretical vs. Practical}:
\begin{itemize}
\item Mathematical infinity vs. finite resources
\item Unknown space requirements handled by allocation
\item Useful computation always finite
\end{itemize}
\end{column}
\begin{column}{0.4\textwidth}
\begin{center}
\textbf{Digital Representation}
\vspace{0.3cm}

\tikz{
% Memory
\draw[thick] (0,3) rectangle (1.5,2);
\node at (0.75,2.5) {RAM};

% Counter
\draw[thick] (0,1.5) rectangle (1.5,1);
\node at (0.75,1.25) {Counter};

% FA
\draw[thick] (0,0.5) rectangle (1.5,0);
\node at (0.75,0.25) {FA};

% Connections
\draw[->] (0.75,1.5) -- (0.75,2);
\draw[->] (0.75,1) -- (0.75,0.5);
}
\end{center}
\end{column}
\end{columns}
\end{frame}

% Section 4: Von Neumann Architecture
\section{Von Neumann Architecture}

\begin{frame}{The Stored Program Concept}
\begin{columns}
\begin{column}{0.6\textwidth}
\concept{Revolutionary Idea}:
\begin{itemize}
\item Programs stored in same memory as data
\item Instructions treated as data
\item Dynamic program modification possible
\item Universal computing machine
\end{itemize}

\vspace{0.3cm}
\concept{Key Principles}:
\begin{itemize}
\item Single memory space for instructions and data
\item Sequential instruction execution
\item Binary representation of information
\item Programmable control unit
\end{itemize}

\vspace{0.3cm}
\concept{Advantages}:
\begin{itemize}
\item Flexibility in programming
\item Self-modifying code capability
\item Simplified hardware design
\item Universal computation model
\end{itemize}
\end{column}
\begin{column}{0.4\textwidth}
\begin{center}
\textbf{Stored Program Concept}
\vspace{0.3cm}

\tikz{
% Memory
\draw[thick] (0,0) rectangle (2,3);
\node at (1,2.7) {Memory};

% Instructions section
\draw[dashed] (0,2) -- (2,2);
\node at (1,2.3) {Instructions};
\node at (1,1.5) {Data};

% CPU
\draw[thick] (-1.5,1) rectangle (-0.5,2);
\node at (-1,1.5) {CPU};

% Connection
\draw[<->] (-0.5,1.5) -- (0,1.5);
}
\end{center}
\end{column}
\end{columns}
\end{frame}

\begin{frame}{Von Neumann Architecture Components}
\begin{columns}
\begin{column}{0.6\textwidth}
\concept{Essential Components}:
\begin{enumerate}
\item \textbf{Central Processing Unit (CPU)}
\item \textbf{Memory Unit}
\item \textbf{Input/Output Devices}
\item \textbf{Control Unit}
\item \textbf{Arithmetic Logic Unit (ALU)}
\end{enumerate}

\vspace{0.3cm}
\concept{CPU Subcomponents}:
\begin{itemize}
\item Control Unit: Instruction fetch and decode
\item ALU: Arithmetic and logical operations
\item Registers: Fast temporary storage
\item Program Counter: Instruction address tracking
\end{itemize}

\vspace{0.3cm}
\concept{Data Flow}:
\begin{itemize}
\item Instructions fetched from memory
\item Data loaded from memory to registers
\item Operations performed in ALU
\item Results stored back to memory
\end{itemize}
\end{column}
\begin{column}{0.4\textwidth}
\begin{center}
\textbf{Von Neumann Architecture}
\vspace{0.3cm}

\tikz{
% CPU
\draw[thick] (0,2) rectangle (2,3);
\node at (1,2.5) {CPU};

% Memory
\draw[thick] (0,0.5) rectangle (2,1.5);
\node at (1,1) {Memory};

% I/O
\draw[thick] (-1,1.25) rectangle (0,1.75);
\node at (-0.5,1.5) {I/O};

% Connections
\draw[<->] (1,2) -- (1,1.5);
\draw[<->] (0,1.5) -- (-1,1.5);
}
\end{center}
\end{column}
\end{columns}
\end{frame}

\begin{frame}{Memory Organization and Addressing}
\begin{columns}
\begin{column}{0.6\textwidth}
\concept{Memory Structure}:
\begin{itemize}
\item Linear array of storage locations
\item Each location has unique address
\item Fixed word size (typically 8, 16, 32, 64 bits)
\item Random access capability
\end{itemize}

\vspace{0.3cm}
\concept{Addressing Schemes}:
\begin{itemize}
\item Direct addressing: Address specifies location
\item Indirect addressing: Address points to address
\item Indexed addressing: Base + offset calculation
\item Relative addressing: PC + displacement
\end{itemize}

\vspace{0.3cm}
\concept{Memory Hierarchy}:
\begin{itemize}
\item Registers: Fastest, smallest capacity
\item Cache: Fast, medium capacity
\item Main memory: Moderate speed, large capacity
\item Secondary storage: Slow, very large capacity
\end{itemize}
\end{column}
\begin{column}{0.4\textwidth}
\begin{center}
\textbf{Memory Organization}
\vspace{0.3cm}

\tikz{
% Memory array
\foreach \y in {0,0.3,0.6,0.9,1.2,1.5} {
  \draw[thick] (0,\y) rectangle (1.5,\y+0.25);
}
\node at (-0.3,1.625) {0};
\node at (-0.3,1.325) {1};
\node at (-0.3,1.025) {2};
\node at (-0.3,0.725) {...};
\node at (-0.3,0.425) {n-1};
\node at (-0.3,0.125) {n};

\node at (0.75,1.625) {Data};
\node at (0.75,1.325) {Instr};
\node at (0.75,1.025) {Data};

% Address bus
\draw[->] (-1,0.8) -- (0,0.8);
\node at (-1.5,0.8) {Address};
}
\end{center}
\end{column}
\end{columns}
\end{frame}

% Section 5: Instruction Set Architecture
\section{Instruction Set Architecture}

\begin{frame}{Instruction Format and Types}
\begin{columns}
\begin{column}{0.6\textwidth}
\concept{Instruction Components}:
\begin{itemize}
\item \textbf{Opcode}: Operation to be performed
\item \textbf{Operands}: Data or addresses for operation
\item \textbf{Addressing Mode}: How operands are specified
\end{itemize}

\vspace{0.3cm}
\concept{Instruction Types}:
\begin{itemize}
\item Data transfer: Load, store, move
\item Arithmetic: Add, subtract, multiply, divide
\item Logical: AND, OR, NOT, XOR
\item Control: Branch, jump, call, return
\item I/O: Input, output operations
\end{itemize}

\vspace{0.3cm}
\concept{Instruction Formats}:
\begin{itemize}
\item Fixed length: All instructions same size
\item Variable length: Instructions vary in size
\item Trade-offs: Simplicity vs. code density
\end{itemize}
\end{column}
\begin{column}{0.4\textwidth}
\begin{center}
\textbf{Instruction Format}
\vspace{0.3cm}

\tikz{
% Instruction format
\draw[thick] (0,2) rectangle (3,2.5);
\draw[thick] (0.8,2) -- (0.8,2.5);
\draw[thick] (1.6,2) -- (1.6,2.5);
\draw[thick] (2.4,2) -- (2.4,2.5);

\node at (0.4,2.25) {Op};
\node at (1.2,2.25) {R1};
\node at (2,2.25) {R2};
\node at (2.7,2.25) {R3};

\node at (0.4,1.7) {6 bits};
\node at (1.2,1.7) {5 bits};
\node at (2,1.7) {5 bits};
\node at (2.7,1.7) {5 bits};

% Example
\node at (1.5,1.2) {ADD R1, R2, R3};
\node at (1.5,0.9) {R1 ← R2 + R3};
}
\end{center}
\end{column}
\end{columns}
\end{frame}

\begin{frame}{Fetch-Decode-Execute Cycle}
\begin{columns}
\begin{column}{0.6\textwidth}
\concept{The Basic Cycle}:
\begin{enumerate}
\item \textbf{Fetch}: Get instruction from memory
\item \textbf{Decode}: Interpret instruction format
\item \textbf{Execute}: Perform the operation
\item \textbf{Store}: Write results back
\end{enumerate}

\vspace{0.3cm}
\concept{Detailed Steps}:
\begin{itemize}
\item PC contains address of next instruction
\item Instruction fetched from memory[PC]
\item PC incremented to next instruction
\item Control unit decodes instruction
\item Operands fetched if needed
\item ALU performs operation
\item Results stored in destination
\end{itemize}

\vspace{0.3cm}
\concept{Control Flow}:
\begin{itemize}
\item Sequential execution by default
\item Branch instructions modify PC
\item Conditional branches test flags
\item Subroutine calls save return address
\end{itemize}
\end{column}
\begin{column}{0.4\textwidth}
\begin{center}
\textbf{Fetch-Decode-Execute}
\vspace{0.3cm}

\tikz{
\node[draw, rectangle, minimum width=2cm, minimum height=0.8cm] (fetch) at (0,3) {Fetch};
\node[draw, rectangle, minimum width=2cm, minimum height=0.8cm] (decode) at (0,2) {Decode};
\node[draw, rectangle, minimum width=2cm, minimum height=0.8cm] (execute) at (0,1) {Execute};
\node[draw, rectangle, minimum width=2cm, minimum height=0.8cm] (store) at (0,0) {Store};

\draw[->] (fetch) -- (decode);
\draw[->] (decode) -- (execute);
\draw[->] (execute) -- (store);
\draw[->] (store) -- ++(1,0) -- ++(0,3.5) -- (fetch);
}
\end{center}
\end{column}
\end{columns}
\end{frame}

% Section 6: Hardware-Software Interface
\section{Hardware-Software Interface}

\begin{frame}{The Architecture Abstraction}
\begin{columns}
\begin{column}{0.6\textwidth}
\concept{Abstraction Layers}:
\begin{itemize}
\item Application software
\item System software (OS, compilers)
\item Instruction Set Architecture (ISA)
\item Microarchitecture
\item Digital logic circuits
\item Analog circuits and devices
\end{itemize}

\vspace{0.3cm}
\concept{ISA as Interface}:
\begin{itemize}
\item Contract between hardware and software
\item Defines visible processor state
\item Specifies instruction semantics
\item Enables software portability
\end{itemize}

\vspace{0.3cm}
\concept{Implementation Independence}:
\begin{itemize}
\item Same ISA, different implementations
\item Performance vs. cost trade-offs
\item Backward compatibility preservation
\end{itemize}
\end{column}
\begin{column}{0.4\textwidth}
\begin{center}
\textbf{Abstraction Layers}
\vspace{0.3cm}

\tikz{
\foreach \y/\label in {0/Devices, 0.5/Analog, 1/Digital, 1.5/Micro-arch, 2/ISA, 2.5/System SW, 3/Application} {
  \draw[thick] (0,\y) rectangle (2,\y+0.4);
  \node at (1,\y+0.2) {\label};
}

% Interface highlight
\draw[thick, red] (-0.1,1.9) rectangle (2.1,2.1);
\node[red] at (2.5,2) {Interface};
}
\end{center}
\end{column}
\end{columns}
\end{frame}

\begin{frame}{Performance Considerations}
\begin{columns}
\begin{column}{0.6\textwidth}
\concept{Performance Metrics}:
\begin{itemize}
\item Execution time (latency)
\item Throughput (bandwidth)
\item Instructions per second (IPS)
\item Clock frequency
\end{itemize}

\vspace{0.3cm}
\concept{Performance Factors}:
\begin{itemize}
\item Instruction count
\item Clock cycles per instruction (CPI)
\item Clock cycle time
\item Memory access time
\end{itemize}

\vspace{0.3cm}
\concept{Performance Equation}:
\begin{align}
\text{CPU Time} &= \text{Instruction Count} \times \text{CPI} \times \text{Clock Cycle Time}
\end{align}

\vspace{0.3cm}
\concept{Optimization Strategies}:
\begin{itemize}
\item Reduce instruction count
\item Improve CPI through pipelining
\item Increase clock frequency
\item Optimize memory hierarchy
\end{itemize}
\end{column}
\begin{column}{0.4\textwidth}
\begin{center}
\textbf{Performance Triangle}
\vspace{0.3cm}

\tikz{
\node at (0,2) {Instruction};
\node at (0,1.7) {Count};
\node at (-1.5,0) {CPI};
\node at (1.5,0) {Clock};
\node at (1.5,-0.3) {Cycle};

\draw[thick] (0,1.5) -- (-1.2,0.3);
\draw[thick] (0,1.5) -- (1.2,0.3);
\draw[thick] (-1.2,0.3) -- (1.2,0.3);

\node at (0,0.8) {CPU Time};
}
\end{center}
\end{column}
\end{columns}
\end{frame}

% Summary
\section{Summary}

\begin{frame}{Chapter Summary}
\begin{columns}
\begin{column}{0.6\textwidth}
\concept{Key Concepts Covered}:
\begin{itemize}
\item Fourth-order systems mark the convergence of physical and informational structures
\item The decision problem led to fundamental computability theory
\item Turing machines provide universal computation model
\item Von Neumann architecture enables stored program concept
\item ISA defines hardware-software interface
\item Performance depends on instruction count, CPI, and clock cycle time
\end{itemize}

\vspace{0.3cm}
\concept{Historical Significance}:
\begin{itemize}
\item From ancient paradoxes to modern computers
\item Mathematical foundations enable practical implementation
\item Architecture abstraction enables software development
\end{itemize}
\end{column}
\begin{column}{0.4\textwidth}
\begin{center}
\textbf{Von Neumann Legacy}
\vspace{0.3cm}

\tikz{
\node[draw, circle, minimum width=2cm] (center) at (0,0) {Von Neumann};
\node[draw, rectangle] (stored) at (0,1.5) {Stored Program};
\node[draw, rectangle] (isa) at (-1.5,0.5) {ISA};
\node[draw, rectangle] (perf) at (1.5,0.5) {Performance};
\node[draw, rectangle] (arch) at (0,-1.5) {Architecture};

\draw[->] (center) -- (stored);
\draw[->] (center) -- (isa);
\draw[->] (center) -- (perf);
\draw[->] (center) -- (arch);
}
\end{center}
\end{column}
\end{columns}
\end{frame}

\begin{frame}{Discussion Questions}
\begin{enumerate}
\item How does the stored program concept enable universal computation?
\item What are the trade-offs between different instruction set architectures?
\item How do abstraction layers facilitate computer system design?
\item What role does the ISA play in software portability?
\item How can performance be optimized at the architecture level?
\item What are the implications of Gödel's incompleteness theorem for computing?
\item How does the Von Neumann bottleneck affect system performance?
\item What alternatives to Von Neumann architecture exist?
\end{enumerate}
\end{frame}

\end{document}

